\chapter{Frequency Response, Bode, and the s-Plane}
\label{ch:splane}
\label{ch:bode}
\Cref{ch:linsys} built first- and second-order systems in the time and
state-space view. This chapter is their frequency and complex-plane companion:
how the same systems respond to steady sinusoids (the Bode plot), and where their
behavior lives in the complex plane (poles, zeros, and the root locus). These are
the two classical design tools---Bode and root locus---and both are just the
transfer function \(G(s)\) read in a particular way.

\section{Decibels: Power, Amplitude, and Roll-Off}
\label{sec:decibels}
Before any plot, we need its vertical axis. A \textbf{decibel} is a logarithmic
measure of a \emph{ratio}, never of an absolute quantity. It is defined on
\emph{power}:
\[
\text{ratio in dB}=10\log_{10}\frac{P_2}{P_1}.
\]
The factor of \(10\) is the whole definition; everything else follows. Because
power in a fixed resistance goes as the square of amplitude, \(P\propto V^2\), an
amplitude ratio picks up a factor of two from the logarithm of a square:
\[
10\log_{10}\frac{P_2}{P_1}
=10\log_{10}\left(\frac{V_2}{V_1}\right)^{\!2}
=20\log_{10}\frac{V_2}{V_1}.
\]
So \(10\log\) for power and \(20\log\) for voltage or current---not two conventions
but one, applied to quantities that differ by a square. The \(20\log\) form is
valid only when the two amplitudes are measured across the \emph{same} impedance;
otherwise return to power.

\begin{workedbox}[: the \(-3\,\mathrm{dB}\) point]
The single most-quoted number in filtering is the \emph{half-power} point. Setting
\(P_2/P_1=\tfrac12\),
\[
10\log_{10}\tfrac12=-3.0103\ldots\approx-\SI{3.01}{\decibel},
\]
which is where the ubiquitous ``\(-\SI{3}{\decibel}\)'' comes from. In amplitude
terms half power means \(V_2/V_1=1/\sqrt2\approx0.707\), and indeed
\[
20\log_{10}\frac{1}{\sqrt2}=-\SI{3.01}{\decibel}
\]
--- the same number, as it must be. This is why the \(-\SI{3}{\decibel}\)
frequency, the half-power frequency, and the \(0.707\)-amplitude frequency are
three names for one point, and why \(-\SI{3}{\decibel}\) is emphatically \emph{not}
half amplitude (half amplitude is \(-\SI{6}{\decibel}\)).
\end{workedbox}

Two frequency-ratio words complete the vocabulary. A \textbf{decade} is a
\(\times10\) frequency ratio; an \textbf{octave} is a \(\times2\) ratio. Since
\(\log_{10}2\approx0.301\), a slope of \(\SI{20}{\decibel}\) per decade is
\(20\times0.301\approx\SI{6}{\decibel}\) per octave---the same slope in different
units.

Finally, \textbf{roll-off} is the term for how fast a response falls once it is
past its passband, quoted as that slope: ``rolls off at
\(-\SI{20}{\decibel}\)/decade'' means each tenfold increase in frequency cuts the
amplitude by a factor of ten. Steeper roll-off means a filter that rejects
out-of-band signals harder.

\section{What a Bode Plot Shows}
A Bode plot is a pair of graphs against logarithmic frequency,
\[
20\log_{10}|G(\jj\omega)|
\quad\text{(magnitude, in dB)},
\qquad
\angle G(\jj\omega)
\quad\text{(phase)}.
\]
The magnitude curve shows gain or attenuation; the phase curve shows lead or lag.
Frequency is logarithmic, so equal ratios take equal space, using the decade and
octave of \cref{sec:decibels}. The frequency response is the transfer function of
\cref{ch:linsys} evaluated on the imaginary axis, \(s=\jj\omega\)---that is, the
system's steady response to a \emph{sinusoidal input} at frequency \(\omega\).

\section{First-Order Response}
The first-order low-pass \(G(s)=1/(1+\tau s)\) has
\[
|G|=\frac{1}{\sqrt{1+(\omega\tau)^2}},
\qquad
\angle G=-\tan^{-1}(\omega\tau).
\]
The magnitude is flat until the \emph{corner frequency} \(\omega_c=1/\tau\), then
rolls off at \(-\SI{20}{\decibel}\) per decade; the phase slides from \(0^\circ\)
to \(-90^\circ\), passing \(-45^\circ\) at the corner (\cref{fig:bode1}). The
corner is the same \(1/\tau\) that set the time constant in \cref{ch:linsys}.

\subsection{Deriving the Asymptotes}
\label{sec:asymptotes}
Both straight lines in \cref{fig:bode1} come from taking limits of \(|G|\), and it
is worth doing once because the same two limits generate every Bode sketch.

\emph{Well below the corner} (\(\omega\tau\ll1\)) the \((\omega\tau)^2\) term is
negligible, so
\[
|G|\approx\frac{1}{\sqrt{1+0}}=1
\qquad\Longrightarrow\qquad
20\log_{10}|G|\approx\SI{0}{\decibel},
\]
a horizontal line at the DC gain. \emph{Well above the corner}
(\(\omega\tau\gg1\)) the \(1\) is negligible instead, so
\[
|G|\approx\frac{1}{\omega\tau}
\qquad\Longrightarrow\qquad
20\log_{10}|G|\approx-20\log_{10}(\omega\tau),
\]
which on a logarithmic frequency axis is a straight line falling
\(\SI{20}{\decibel}\) for every tenfold increase in \(\omega\): the
\(-\SI{20}{\decibel}\)/decade roll-off. The two asymptotes meet where
\(1=\omega\tau\), i.e.\ at \(\omega=1/\tau\)---which is what makes \(\omega_c\) the
``corner.''

\emph{Why the exact curve misses the corner by \(\SI{3}{\decibel}\).} The
asymptotes each dropped one term; exactly at the corner neither term is negligible,
because \(\omega\tau=1\) makes them \emph{equal}. There,
\[
|G(\jj\omega_c)|=\frac{1}{\sqrt{1+1}}=\frac{1}{\sqrt2},
\qquad
20\log_{10}\frac{1}{\sqrt2}=-\SI{3.01}{\decibel},
\]
the half-power point of \cref{sec:decibels}. So the \(\SI{3}{\decibel}\) gap
between the asymptote intersection and the true curve is not an error in the sketch
--- it is the exact size of the approximation at the one frequency where both terms
matter equally. The phase behaves the same way: \(-\tan^{-1}(1)=-45^\circ\), exactly
halfway through its \(0^\circ\to-90^\circ\) travel.

\begin{figure}[htbp]
\centering
\includegraphics[width=0.82\textwidth]{bode_first_order_lowpass}
\caption{First-order low-pass Bode plot. The corner is at \(\omega_c=1/\tau\),
where the gain is \(-\SI{3}{\decibel}\) and the phase is \(-45^\circ\).}
\label{fig:bode1}
\end{figure}

\section{Building the Sketch from Poles and Zeros}
A Bode magnitude plot can be sketched by hand from the pole and zero locations,
because each contributes a straight-line asymptote that adds to the others:
\begin{itemize}
\item Start at the low-frequency gain (the DC value \(|G(0)|\)).
\item At each \textbf{pole} corner, bend the slope \emph{down} by
\(\SI{20}{\decibel}\) per decade and subtract up to \(90^\circ\) of phase.
\item At each \textbf{zero} corner, bend the slope \emph{up} by
\(\SI{20}{\decibel}\) per decade and add up to \(90^\circ\) of phase.
\item A corner sits at the frequency equal to that pole's or zero's distance from
the origin; an \(n\)-pole low-pass therefore rolls off at about
\(-20n\,\si{\decibel}\) per decade far above cutoff
(\cref{sec:cascadeadd} derives that rule from the fact that cascaded sections add
in decibels).
\end{itemize}
The exact curve rounds the corners---down \(\SI{3}{\decibel}\) at a first-order
pole, and peaking up at a lightly damped second-order pair (next section)---but
the asymptotes get the shape right.

\section{Second-Order Response: the Resonant Peak}
Recall the standard second-order transfer function from \cref{ch:linsys},
\[
G(s)=\frac{\omega_0^2}{s^2+2\zeta\omega_0 s+\omega_0^2},
\qquad
\text{poles at } s=-\zeta\omega_0\pm\jj\omega_0\sqrt{1-\zeta^2},
\]
and evaluate it on the imaginary axis, \(s=\jj\omega\):
\[
G(\jj\omega)=\frac{\omega_0^2}{(\omega_0^2-\omega^2)+\jj\,2\zeta\omega_0\omega},
\qquad
|G(\jj\omega)|=\frac{\omega_0^2}
{\sqrt{(\omega_0^2-\omega^2)^2+(2\zeta\omega_0\omega)^2}} .
\]
At \(\omega=\omega_0\) the first bracket vanishes and only the damping term
survives, leaving \(|G|=1/(2\zeta)=Q\): the peak height \emph{is} the quality
factor. Geometrically, sampling \(G(\jj\omega)\) near \(\omega_0\) brings the
evaluation point close to the pole pair, producing a resonant peak whose height
grows with \(Q\) (\cref{fig:bode2}). The flattest response with no peak occurs at
\(Q=1/\sqrt{2}\approx0.707\) (the Butterworth value); above it the magnitude
peaks near \(\omega_0\) by a factor of about \(Q\), and the peak sharpens as the
poles approach the \(\jj\omega\) axis.

\begin{figure}[htbp]
\centering
\includegraphics[width=0.82\textwidth]{bode_second_order}
\caption{Second-order magnitude response. Higher \(Q\) (lower \(\zeta\)) means a
taller, sharper resonant peak, because the poles sit closer to the \(\jj\omega\)
axis.}
\label{fig:bode2}
\end{figure}

\section{Reading a System off the Plot}
Given a measured or computed Bode plot, four features tell you what a
second-order LTI system (for example, a circuit) does:
\begin{description}
\item[Passband and cutoff] the flat region is the passband; the
\(-\SI{3}{\decibel}\) frequency (magnitude \(1/\sqrt2\)) is the cutoff or corner.
\item[Roll-off slope] the far-out slope gives the order: \(-20n\,\si{\decibel}\)
per decade means \(n\) poles. Steeper skirts reject out-of-band signals harder.
\item[Resonant peak] a bump near \(\omega_0\) signals a high-\(Q\) pole pair; its
height grows with \(Q\) and its \(-\SI{3}{\decibel}\) width is the bandwidth
\(BW\approx f_0/Q\).
\item[Phase] the phase slope near a corner flags how much lead or lag the stage
adds---the quantity that governs stability once the stage sits inside a feedback
loop (\cref{ch:feedback}).
\end{description}

\section{The s-Plane: Poles, Zeros, and Stability}
The Bode plot samples \(G(s)\) only on the imaginary axis---that is, only for
sinusoidal inputs. The s-plane shows the whole picture: where \(G\) blows up
(poles) and vanishes (zeros) across the entire complex plane
\(s=\sigma+\jj\omega\), growth and decay (\(\sigma\)) included.

\begin{controlsbox}
\textbf{Refresher: poles and zeros.} Write \(G(s)=N(s)/D(s)\). The \emph{poles}
are the roots of \(D\) (the natural modes, equal to the eigenvalues of \(A\)); the
\emph{zeros} are the roots of \(N\) (frequencies transmitted as zero). An unforced
mode evolves as \(\ee^{s_i t}\) with \(s_i=\sigma_i+\jj\omega_i\), so the sign of
the real part is the entire stability story: left half-plane (\(\sigma_i<0\))
decays and is stable; the \(\jj\omega\) axis sustains a pure oscillation (the
knife-edge an oscillator sits on); the right half-plane grows and is unstable.
Zeros do not affect stability, but they shape how the stable modes appear.
\end{controlsbox}

\subsection{Reading the Pole Geometry}
For the second-order pole pair \(s=-\zeta\omega_0\pm\jj\omega_d\), two facts make
the plot readable by eye (\cref{fig:poles2}). First, while underdamped the poles
lie on a circle of radius \(\omega_0\), because
\(|s|^2=(\zeta\omega_0)^2+(\omega_0\sqrt{1-\zeta^2})^2=\omega_0^2\); the distance
from the origin is the natural frequency. Second, the angle \(\theta\) from the
negative real axis satisfies \(\cos\theta=\zeta\), so the angle encodes the
damping and therefore \(Q\). The real part \(\alpha=\zeta\omega_0\) is the
envelope decay rate \emph{and} half the half-power bandwidth: the full bandwidth is
\[
BW=2\alpha=\frac{\omega_0}{Q}\quad[\si{\radian\per\second}].
\]
\emph{Distance from the imaginary axis is half the bandwidth.} A pole hugging the
axis rings for a long time and shows a narrow, tall peak; a pole far from the axis
settles quickly and shows a broad, low peak.

\begin{figure}[htbp]
\centering
\includegraphics[width=0.56\textwidth]{poles_second_order}
\caption{Second-order poles for several damping ratios. They ride a circle of
radius \(\omega_0\); the angle from the negative real axis has
\(\cos\theta=\zeta\). \(\zeta=0\) sits on the \(\jj\omega\) axis (undamped).}
\label{fig:poles2}
\end{figure}

\subsection{Zeros and the Pole-Zero Plot}
A pole-zero plot marks poles with \(\times\) and zeros with \(\circ\)
(\cref{fig:pz-notch}). Poles set the resonance and damping; a zero on or near the
\(\jj\omega\) axis pulls the magnitude \emph{down} toward a null at that
frequency---the mechanism of a notch, trap, or band-stop.

\begin{figure}[htbp]
\centering
\includegraphics[width=0.52\textwidth]{poles_zeros_notch}
\caption{Pole-zero plot of a notch response. The complex pole pair
(\(\times\)) just inside the left half-plane sets the resonance; the transmission
zeros (\(\circ\)) on the \(\jj\omega\) axis force the output to zero at the
rejection frequency \(\omega_0\).}
\label{fig:pz-notch}
\end{figure}

\section{Root Locus}
When a parameter varies---a resistance, a gain, a feedback fraction---the poles
move, tracing a \emph{root locus}. For the canonical second-order system
\(s^2+2\zeta\omega_0 s+\omega_0^2\), sweeping the damping \(\zeta\) from zero
upward traces the path in \cref{fig:rootlocus2}: the poles begin on the
\(\jj\omega\) axis at \(\pm\jj\omega_0\) (undamped), arc along the circle of
radius \(\omega_0\) into the left half-plane, meet at \(-\omega_0\) at critical
damping (\(\zeta=1\)), then split and travel apart along the negative real axis
(overdamped).

\begin{figure}[htbp]
\centering
\includegraphics[width=0.58\textwidth]{rootlocus_second_order}
\caption{Root locus of a second-order system as damping increases. Poles begin
on the \(\jj\omega\) axis (undamped), ride the \(\omega_0\) circle to \(-\omega_0\)
at critical damping, then separate along the real axis (overdamped).}
\label{fig:rootlocus2}
\end{figure}

The series RLC circuit of \cref{ch:rlc} is exactly this locus with
\(\zeta=\tfrac{R}{2}\sqrt{C/L}\), so sweeping \(R\) walks the poles from a pure
LC oscillation (\(R=0\)) to a heavily damped pair; \cref{ch:rlc} draws that
specific locus.

\section{Worked Example: Reading a High-Q Pole Pair}
Take a lightly damped second-order system with natural frequency
\(\omega_0=\SI{2.13e7}{\radian\per\second}\) and damping ratio \(\zeta=0.0117\)
(equivalently \(Q=1/(2\zeta)\approx42.6\)). Its poles are
\[
s=-\zeta\omega_0\pm\jj\,\omega_0\sqrt{1-\zeta^2}
\approx-2.5\times10^{5}\pm\jj\,2.13\times10^{7}\ \si{\radian\per\second}.
\]
The real part is about \(1/85\) of the imaginary part, so the poles are pressed
against the \(\jj\omega\) axis: a very high-\(Q\), sharply peaked response that
rings for roughly \(Q\approx43\) cycles before decaying. Nothing here is
electrical---these are just numbers read off the standard second-order form.
\Cref{ch:rlc} realizes exactly this pole pair as a tuned LC circuit, showing
which components \(L\), \(C\), and \(R\) put the poles here.

\begin{exambox}
The exam does not ask you to plot poles or draw Bode diagrams. It asks about
cutoff, roll-off slope, resonant frequency, \(Q\), bandwidth, selectivity,
ringing, oscillator stability, filter sharpness, and traps---all of which answer
one question: \emph{where are the poles and zeros?} Near the \(\jj\omega\) axis
means sharp and selective; on the axis means sustained oscillation; a zero on the
axis means a rejected frequency.
\end{exambox}
